% ==============================================================================
% ALCANCE
% ==============================================================================
% Contenido provisto directamente por el estudiante. Dos correcciones
% aplicadas antes de incorporarlo:
%   1. "la sección 1.3" (referencia manual, en Alcance operativo) se
%      reemplazó por \label/\ref sobre la subsección "Alcance operativo":
%      con Alcance como capítulo propio, esa subsección será algo como
%      "5.1.3" según dónde quede el capítulo, no "1.3" (número que parece
%      arrastrado de una numeración distinta, probablemente la del
%      protocolo original).
%   2. "la sección de factibilidad de la metodología de este protocolo" se
%      corrigió a "de este documento": este texto ya es parte de la tesis,
%      no del protocolo.
%
% NOTA ESTRUCTURAL (sin resolver): "Alcance" no aparece en la lista de 10
% capítulos de .claude/CLAUDE.md (Introducción...Recomendaciones). Sí existe
% como capítulo opcional de la plantilla UVG (\CAPalcance). Se incluye por
% instrucción explícita del estudiante; el índice final tendrá 11 capítulos
% en vez de 10 si esto se mantiene así.
%
% Corrección 2026-10-02 (criterio de latencia del middleware completo, en
% "Alcance técnico"): el criterio original fijaba un incremento relativo de
% p95 del 5% bajo carga normal y 10% bajo carga de ataque simulado. Diego
% Valenzuela (componente de IA del mismo proyecto integrador) señaló por
% chat (2026-10-02) que un criterio relativo es débil cuando la línea base
% de latencia del servidor es de pocos milisegundos: un 5% de esa base es
% una fracción de milisegundo, difícil de medir de forma confiable y poco
% significativa como criterio de aceptación. Reportó que a él le funcionó
% un límite absoluto de 15~ms en su propio componente, y recomendó usar un
% número fijo en ms en vez de un porcentaje salvo que la línea base ya sea
% grande (> 100~ms). El estudiante confirmó adoptar un único umbral
% absoluto de 15~ms (sin distinguir carga normal de carga de ataque
% simulado) para el middleware completo, reemplazando el 5%/10% relativo.
% El criterio de 5~ms del pipeline de extracción de características
% (aislado, dentro del worker thread) no cambia: ya era un valor absoluto.
\ifdefined\modoensamble
\else
	\input{0-preambulo_standalone}
\fi

\chapter{Alcance}

\section{Alcance del proyecto}

\subsection{Alcance funcional}
El proyecto comprende el diseño, implementación y evaluación de un middleware distribuible
como paquete npm para aplicaciones Node.js basadas en Express, capaz de detectar y responder,
en tiempo de ejecución, a peticiones HTTP correspondientes a cuatro vectores de ataque
intra-request: SQL Injection, Cross-Site Scripting (XSS), Path Traversal/LFI e inyección de
comandos. La respuesta del sistema ante una petición clasificada como maliciosa se limita a dos
mecanismos verificables: el bloqueo de la petición mediante código de respuesta HTTP 403 y el
registro de una alerta. El proyecto entrega, como producto funcional, el paquete npm publicado, los
scripts de construcción del corpus etiquetado y los vectores de características derivados, y la
documentación técnica correspondiente. No se entregan los datos crudos de las fuentes públicas
empleadas: cada fuente se obtiene de su origen conforme a su licencia.

\subsection{Alcance técnico}
\label{subsec:alcance-tecnico}
El sistema se integra exclusivamente como middleware dentro del pipeline de peticiones de
Express.js, sin requerir modificaciones al runtime de Node.js ni infraestructura de red externa
(proxy inverso o plataforma SaaS de terceros). El componente de inferencia consume un modelo
serializado en formato ONNX mediante \texttt{onnxruntime-node}, ejecutado dentro de un
\texttt{worker\_thread} independiente para evitar el bloqueo del Event Loop principal. El alcance
técnico está delimitado por dos restricciones cuantitativas, definidas como criterios de aceptación
verificables: el pipeline de extracción de características debe ejecutarse en un tiempo igual o
menor a 5~ms por petición, medido de forma aislada dentro de su worker thread; y el middleware
completo no debe incrementar la latencia p95 del servidor anfitrión en más de 15~ms absolutos
respecto de la línea base sin middleware activo, tanto bajo carga normal como bajo carga de
ataque simulado. Se adopta un umbral absoluto en lugar de uno relativo porque, en un servidor
cuya línea base de latencia es de pocos milisegundos, un porcentaje fijo representaría una
fracción de milisegundo: un criterio estadísticamente frágil y difícil de verificar de forma
confiable bajo las condiciones de medición de este proyecto.

\subsection{Alcance operativo}
\label{subsec:alcance-operativo}
La validación principal del sistema se realiza en un entorno de simulación controlado y
reproducible, construido mediante Docker Compose, con generación de tráfico legítimo mediante
Artillery y de tráfico malicioso mediante herramientas especializadas por vector de ataque
(sqlmap, OWASP ZAP y scripts propios para Path Traversal/LFI y Command Injection). La
validación extendida con entornos productivos de startups guatemaltecas, contemplada como
escenario complementario en la Fase 4 del cronograma, queda condicionada a la disponibilidad de
un acuerdo de colaboración con dichas organizaciones y no constituye el criterio de aceptación
principal del proyecto.

\section{Delimitación}

\subsection{Lo que comprende este documento}
Este documento delimita específicamente el componente de ciberseguridad del sistema
\textit{logSguarDian}: el análisis y selección de vectores de ataque, la construcción del entorno de
simulación y el dataset etiquetado, la implementación inicial del extractor de características y el paso a \texttt{worker\_threads}, la capa de captura de peticiones HTTP y el sistema de
respuesta en tiempo de ejecución, y la evaluación de cobertura de detección mediante pruebas de
penetración controlada. La integración de estos componentes con el modelo de aprendizaje
automático, así como las pruebas de integración extremo a extremo y la documentación final
conjunta, son responsabilidad compartida de ambos integrantes del proyecto y se documentan en
la sección de metodología correspondiente.

\subsection{Lo que queda fuera del alcance}
\label{subsec:fuera-alcance}
Explícitamente quedan fuera del alcance de este trabajo:

\begin{itemize}
    \item El diseño, entrenamiento, selección de arquitectura y serialización del modelo de
    aprendizaje automático, documentado de forma independiente en el trabajo de graduación de
    Diego Valenzuela, correspondiente al componente de Inteligencia Artificial.
    \item La detección de vectores de ataque distintos a los cuatro delimitados en el Objetivo
    Específico 1 (SQL Injection, XSS, Path Traversal/LFI e inyección de comandos). Categorías
    como Server-Side Request Forgery (SSRF), Cross-Site Request Forgery (CSRF) o ataques de
    denegación de servicio a nivel de aplicación no forman parte del sistema de detección
    propuesto.
    \item El soporte para runtimes distintos a Node.js o frameworks distintos a Express.js. El
    sistema no está diseñado para operar sobre aplicaciones construidas en Django, Spring,
    Laravel u otros stacks tecnológicos.
    \item El reemplazo de mecanismos perimetrales de seguridad existentes. El middleware se
    plantea como una capa complementaria a nivel de aplicación, no como sustituto de un
    Web Application Firewall, un IDS de red o un SIEM empresarial.
    \item El cumplimiento normativo o certificación bajo cualquier marco legal de ciberseguridad
    vigente o en trámite en Guatemala. El proyecto no garantiza ni evalúa conformidad legal;
    su relevancia frente al panorama regulatorio se documenta como antecedente contextual,
    no como criterio de aceptación.
    \item El soporte, mantenimiento o actualización del paquete npm posterior a la fecha de entrega del informe final en noviembre de 2026.
    \item El despliegue en infraestructura de producción de terceros fuera del entorno de
    validación controlado y del escenario complementario descrito en la Sección~\ref{subsec:alcance-operativo},
    cuya realización depende de un acuerdo de colaboración externo al proyecto.
\end{itemize}

\section{Restricciones y supuestos}
El desarrollo del proyecto está sujeto a las siguientes restricciones y supuestos, verificables al
momento de redacción de este documento:

\begin{itemize}
    \item \textbf{Restricción temporal.} El proyecto se ejecuta entre mayo y octubre de 2026,
    organizado en seis fases con dependencias secuenciales explícitas; la Fase 3, correspondiente
    al desarrollo del middleware, depende de la finalización de la Fase 1 (dataset) y la Fase 2
    (entrenamiento del modelo, a cargo del componente de IA).
    \item \textbf{Restricción presupuestaria.} El proyecto opera con un presupuesto de USD 0 en
    condiciones normales, con un techo de USD 50 en el escenario de despliegue cloud temporal
    para validación externa. Esta restricción excluye explícitamente el uso de plataformas de
    seguridad comerciales con licenciamiento de pago como parte del sistema desarrollado.
    \item \textbf{Restricción de recursos humanos.} El componente de ciberseguridad es
    responsabilidad de un único integrante, con dedicación de 20 horas semanales durante las
    Fases 1 a 3 y 15 horas semanales durante las Fases 4 a 6, sin personal externo adicional
    disponible.
    \item \textbf{Supuesto de disponibilidad de dependencias.} Se asume la disponibilidad
    continua de las herramientas de código abierto empleadas (Node.js, Express, Docker,
    sqlmap, OWASP ZAP, Jest) y del dataset público de referencia (OWASP ModSecurity Dataset,
    Zenodo) durante toda la duración del proyecto; el riesgo de interrupción del acceso al
    dataset se mitiga mediante copia local descargada al inicio de la Fase 1.
    \item \textbf{Supuesto de estabilidad del entorno de simulación.} Se asume que el entorno
    Docker Compose reproduce condiciones de tráfico suficientemente representativas del
    comportamiento de una API REST en producción para que los resultados obtenidos en el
    entorno controlado sean generalizables, supuesto que la validación complementaria con
    startups guatemaltecas busca contrastar sin condicionar a ella el cumplimiento del proyecto.
\end{itemize}

\section{Viabilidad del alcance propuesto}
La viabilidad del alcance definido descansa en tres condiciones verificables. Primero, la totalidad
del stack tecnológico requerido (Node.js, Express, onnxruntime-node, Docker) está disponible,
mantenido activamente y documentado para los casos de uso específicos del proyecto, sin
dependencia de tecnología en desarrollo o no probada. Segundo, el dataset base de referencia está
publicado bajo licencia abierta y ha sido empleado previamente en investigaciones de detección de
intrusiones en aplicaciones web, lo que reduce el riesgo de que la construcción del entorno de
simulación se convierta en un cuello de botella del cronograma. Tercero, el alcance funcional se
limita deliberadamente a cuatro vectores de ataque y a dos mecanismos de respuesta, en lugar de
plantear una cobertura exhaustiva de OWASP Top 10, precisamente para mantener el proyecto
ejecutable dentro de las 20 semanas efectivas de desarrollo disponibles y de los recursos humanos
descritos en la sección anterior. Los riesgos identificados que podrían afectar el cumplimiento del
alcance definido, junto con sus planes de mitigación, se documentan de forma exhaustiva en la
sección de factibilidad del Capítulo de Metodología de este documento.

\ifdefined\modoensamble
\else
	\end{document}
\fi
